% sec9_5.tex — Section 9.5: Which Unit-Root Test to Use?

Besides the DF and ADF tests, a number of other unit-root tests are available (see
Maddala and Kim, 1998, Chapter~3, for a catalogue). The most prominent among
them are the Phillips (1987) test for case~(4) (mentioned at the end of the previous
section) and the Phillips-Perron (1988) test, a generalization of the Phillips test
to cover cases~(5) and~(6). (The Phillips tests are derived in Analytical Exercise~6.)
Their tests are based on the OLS estimate of the $y_{t-1}$ coefficient in an AR(1)
equation, rather than an augmented autoregression, with the long-run variance of
$\{\Delta y_t\}$ estimated from the residuals of the AR(1) equation. We did not present
them in the text because the finite-sample properties of the tests are rather poor
(see below). There is a new generation of unit-root tests with reasonably low size
distortions and good power. They include the ADF-GLS test of Elliott, Rothenberg,
and Stock (1996) and the $M$-tests of Perron and Ng (1996).

\subsection*{Local-to-Unity Asymptotics}

These and most other unit-root tests are consistent against all fixed I(0) alternatives.\footnote{That the DF tests are consistent against general I(0) alternatives was verified in Review Question~5 to Section~9.3. For the consistency of the ADF tests, see, e.g., Stock (1994, p.~2770).} Which consistent test should be chosen over the other? Recall that in
Section~2.4 we defined the asymptotic power against a sequence of local alternatives that drift toward the null at rate $\sqrt{T}$. It turns out that in the unit-root context
the appropriate rate is $T$ rather than $\sqrt{T}$ (see, e.g., Stock, 1994, pp.~2772--2773).
That is, the probability of rejecting the null of $\rho = 1$ when the true DGP is
\begin{equation}
\label{eq:9.5.1}
\rho = 1 + \frac{c}{T}
\end{equation}
has a limit between 0 and 1 as $T$ goes to infinity. This limit is the asymptotic
power. The sequence of alternatives described by~\eqref{eq:9.5.1} is called \textbf{local-to-unity
alternatives}. By definition, the asymptotic power equals the nominal size when
$c = 0$. It can be shown that the DF or ADF $\rho$ tests are more powerful than the
DF or ADF $t$ tests against local-to-unity alternatives (see Stock, 1994, pp.~2774--2776). That is, for any nominal size and $c$, the asymptotic power of the DF or ADF
$\rho$ statistic is greater than that of the DF or ADF $t$-statistic. On this score, we should
prefer the $\rho$-based test to the $t$-based test, but the verdict gets overturned when we
examine their finite-sample properties.

\subsection*{Small-Sample Properties}

Turning to the issue of testing an I(1) null against a fixed, rather than drifting, I(0)
alternative, the two desirable finite-sample properties of a test are a reasonably low
size distortion and high power against I(0) alternatives. There is a large body of
Monte Carlo evidence on the finite-sample properties of the DF, ADF, and other
unit-root tests (for a reference, see the papers cited on p.~2777 of Stock, 1994, and
see Ng and Perron, 1998, for a comparison of the ADF-GLS and the $M$-tests).
Major findings are the following.

\begin{itemize}
\item Virtually all tests suffer from size distortions, particularly when the null is in a
sense close to being I(0). The $M$-test generally has the smallest size distortion,
followed by the ADF $t$. The size distortion of the ADF $\rho$ is much larger than
that of the ADF $t$. For example, the I(1) null considered by Schwert (1989) is
\begin{equation}
\label{eq:9.5.2}
y_t - y_{t-1} = \varepsilon_t + \theta \varepsilon_{t-1}, \quad
\{\varepsilon_t\} \text{ Gaussian i.i.d.}
\end{equation}
When $\theta$ is close to $-1$, this I(1) process is very close to a white noise process,
because the AR root of unity nearly equals the MA root of $-\theta$. Simulations by
Schwert (1989) show that the ADF $\rho$ and Phillips-Perron $Z_\rho$ and $Z_t$ tests (see
Analytical Exercise~6) reject the I(1) null far too often in finite samples when $\theta$
is close to $-1$. In particular, for the Phillips-Perron tests, the actual size when
the nominal size is 5 percent is more than 90 percent even for sample sizes as
large as 1,000 when $\theta = -0.8$.

\item In the case of ADF tests, how the lag length $\hat{p}$ is selected affects the finite-sample size and power. For a given rule for choosing $\hat{p}$, the power is generally
higher for the ADF $\rho$-test than for the ADF $t$-test.

\item The power is generally much higher for the ADF-GLS and the $M$-tests than for
the ADF $t$- and $\rho$-tests.
\end{itemize}

Unlike the $M$-statistic, the ADF-GLS statistic can be calculated easily by standard
regression packages. In the empirical exercise of this chapter, you will be asked to
use the ADF-GLS to test an I(1) null.
